\section{Experiments}
\label{sec:experiments}

\subsection{Questions and protocol}

% ROLE: overview -- define the four questions and the common evidence standard.
The experiments ask four questions: (Q1) do canonical coordinates remain stable
when raw latent charts change; (Q2) do compact equations transfer to unseen
future entities; (Q3) does a decoder prior add value specifically when support
does not identify the symbolic structure; and (Q4) what predictive accuracy is
traded for a compact named expression?  We split by entity and, where available,
DOI and time.  A test entity's coefficients and calibrated state use support
targets only; query targets are read once for scoring.  Primary reporting uses
physical-unit pooled $R^2$, entity- or DOI-cluster bootstrap intervals, the
fraction of entities above $R^2=0.85$, and paired comparisons.  The heterogeneous
Crystal stress test additionally reports median/p95/maximum entity NRMSE and
negative-prediction count; every displayed aggregate requires finite full query
coverage.  For fitted methods with an explicit perturbation audit, perturbing
hidden query targets must leave all fitted inputs, coefficients, and predictions
unchanged; all such declared differences are zero.
For entity $e$, NRMSE is query RMSE divided by the population standard deviation
of its query targets ($\mathrm{ddof}=0$).  Controlled five-seed GIRD predictions
are aggregated pointwise before entity metrics; neural-ablation seed rules are
reported with their respective results.
The real-data cohorts come from the Starrydata thermoelectric collection
\citep{katsura2025starrydata} and the NIST ThermoML archive
\citep{riccardi2021thermoml,frenkel2006thermoml}.  Frozen cohort manifests bind
the selected source files by repository-local SHA-256; mutable ``latest'' links
are not treated as the experimental artifact.

% ROLE: fairness -- state information paths and baseline families.
Every method is labeled by its information path.  Support structure re-q fits a
frozen named basis; raw-q ridge maps neural coordinates directly to equation
coefficients; decoder-functional projection probes the calibrated decoder;
GIRD fuses that projection using a validation-selected weight fixed within each
predeclared support regime; rank and conditioning diagnose the regime rather
than route individual entities.  Comparators include
the same-budget direct-target dictionary, FPCA \citep{ramsay2005functional},
CNP, a temperature-only no-q MLP, and support nearest-neighbor, linear, kNN,
and PCHIP predictors \citep{fritsch1984pchip}.  The
strongest support-aware baseline is displayed beside each headline expression
score.  Expression fidelity and predictive superiority are tested separately.

\subsection{Gauge intervention and controlled prior evidence}

% ROLE: evidence -- verify the representation-level invariance claim.
We first intervene on the latent chart while preserving the decoder response.
Across three controlled response families, five training seeds, and 25 exact
affine gauges per family--seed cell (375 interventions), the constructed
counterfactual changes raw coordinates by as much as $4.996$ while changing
decoder predictions and fixed-basis coefficients by at most
$4.50\times10^{-15}$ and $2.04\times10^{-14}$.  This tests the algebraic
response-quotient claim.  A separate 75-case experiment independently
recalibrates each of five affine charts per family--seed cell.  The original
response-metric Gauss--Newton extension diverges by as
much as $0.0448$ in query response; float64 direct least squares, a response-loss
acceptance margin, and
15 fixed Gauss--Newton steps reduce the maximum response and coefficient
differences to $3.66\times10^{-10}$ and $6.23\times10^{-9}$, with identical
paired line-search decisions.  This narrower result supports numerical affine
equivariance only under the stated full-rank calibration assumptions.

% ROLE: boundary evidence -- show that GIRD has a conditional, not universal, mechanism.
The controlled GIRD study next varies support identifiability.  Under four-point
support, aggregate polynomial and thermodynamic selections match their
$\lambda=0$ results.
For the relaxation family under four-point support, a finite prior reduces error
from $0.05277$ to $0.03435$, a $34.9\%$ development improvement.  FPCA remains
better than every GIRD variant, and two of 30 family--seed--regime OMP
certificate audits fail.
These results are consistent with the rank-conditioned bias--variance mechanism
and expose its failure modes; they do not establish a universal prior
advantage.

\subsection{One-shot real-data expression transfer}

% ROLE: primary evidence -- report the two completed temporal confirmations in one table.
Table~\ref{tab:temporal-expression} reports two single-use temporal
confirmations.  ZT's quadratic endpoint was fixed a priori.  Vapor alone used
structure selection: five DOI-grouped folds scored exactly
$\{\log(T/T_r),(T-T_r)/T_r,(1/T-1/T_r)^2\}$ as one-term corrections to coarse
Clausius--Clapeyron by median physical NRMSE over all 282 held-out entities;
candidates within 1\% use the displayed order.  ZT development queries evaluate
its fixed endpoint and vapor development queries select the correction.  After
sealing, future targets only score the structure and support estimates its
coefficients (full audit in Appendix~\ref{app:protocol}).

\begin{table*}[t]
  \centering
  \small
  \caption{One-shot unseen-entity temporal expression results.  Confidence
  intervals resample entities, with an additional DOI-cluster interval for
  vapor pressure.  ``Pass'' counts entities with individual physical
  $R^2\geq0.85$.  The strongest support-aware baseline remains adjacent; higher
  baseline accuracy does not invalidate the separate expression endpoint.
  Here $\tau=(T-280.584300\,\mathrm{K})/(193.386207\,\mathrm{K})$,
  $T_r=362.72\,\mathrm{K}$, and $\Delta T^{-1}=1/T-1/T_r$.}
  \label{tab:temporal-expression}
  \setlength{\tabcolsep}{4pt}
  \begin{tabular}{@{}llrrrrl@{}}
    \toprule
    Domain & Named structure & $N_e$ & $N_q$ & Physical $R^2\uparrow$ & $\geq.85$ & Best support $R^2\uparrow$ \\
    \midrule
    Starry ZT
      & $c_0+c_1\tau+c_2\tau^2$
      & 30 & 919 & \textbf{0.988810} & 20/30 & kNN 0.982678 \\
    Vapor pressure
      & $c_0+c_1\Delta T^{-1}+c_2\log(T/T_r)$
      & 84 & 1,745 & \textbf{0.999581} & 83/84 & log-PCHIP 0.999937 \\
    \bottomrule
  \end{tabular}
\end{table*}

% ROLE: ZT interpretation and limitation -- connect coefficient meaning, stability, and paired baseline evidence.
Here $c_i$ denotes a canonical response coefficient, never a raw neural latent
coordinate.  For ZT, $c_0$, $c_1$, and $c_2$ encode a reference value,
first-order
temperature sensitivity, and curvature.  Entity-bootstrap physical $R^2$ is
$[0.973306,0.994708]$, and median entity $R^2$ is $0.941739$.  On development
entities, changing the stratified-support offset gives response-coordinate
distance Spearman $0.9780$, $0.9733$, and $0.9679$ relative to offset zero; the
physical curvature rank has median/minimum pairwise Spearman
$0.9593/0.9397$.  The expression has lower NRMSE than kNN on 16/30 future
materials, but the paired Wilcoxon test is nonsignificant ($p=0.465$).  Its
scientific test is therefore the predicted slope reversal near
$\tau^\star=-c_1/(2c_2)$, not a claim of universal predictive gain.

% ROLE: vapor-pressure interpretation and limitation -- expose the strongest real symbolic result and coefficient conditioning.
For vapor pressure, the coordinates give reference log pressure $c_0$
(equivalently $P_r=(1\,\mathrm{kPa})e^{c_0}$), effective reference
vaporization enthalpy $R(c_2T_r-c_1)$, and an effective temperature correction
$Rc_2$.  This reading uses an ideal-gas, negligible-liquid-volume,
finite-temperature-range linearization and therefore defines effective response
coordinates, not exact calorimetric quantities.  Entity- and DOI-bootstrap
physical $R^2$ intervals are
$[0.998880,0.999983]$ and $[0.998562,0.999984]$.  Across four development
support offsets, whose fits use outer-training-fold median reference
temperatures rather than the temporal confirmation's frozen $T_r$, the three
named-coordinate rank correlations have medians
$0.9991$, $0.9762$, and $0.7974$; joint response-coordinate geometry has
median/minimum correlation $0.8905/0.8546$.  The weakest correction term is
limited by design collinearity (each support-design column is divided by its
$\ell_2$ norm before the condition number; median/p95/maximum
$135/917/6417$) and is not presented as a calorimetric measurement.  Log-pressure PCHIP has
higher pooled $R^2$, while the compact expression provides a global named
response chart and a falsifiable comparison to independent enthalpy data.

\subsection{What the neural representation contributes}

% ROLE: ablation -- distinguish raw q, decoder response, and direct support fitting.
The neural ablations separate coordinate readability from complete bridge
success.  On the external vapor-pressure cohort, decoder-functional projection
reaches physical $R^2=0.990390$, whereas a direct raw-q ridge expression reaches
only $0.304492$; both are pointwise medians of three seed predictions in physical
pressure space.  The functional projection nevertheless fails its stronger
cross-seed geometry gate and is not claim-eligible as the confirmed expression;
the $0.999581$ result comes from support structure re-q.  The same qualitative
separation occurs for ZT, where each outer-fold prediction is the pointwise
median over three seeds, fixed-basis decoder-functional fidelity is high
but extreme entity-NRMSE tails prevent a universal neural-bridge claim.  These
ablations show that response projection repairs readability of a learned
function, not that every trained decoder yields a safe equation prior.

\subsection{Crystal heat capacity stress test}

% ROLE: development evidence -- report all baselines and tails before opening confirmation.
Crystal heat capacity tests a heterogeneous two-stage response.  Ordinary
curves use NIST's five-term Shomate form \citep{nist2025webbook}; support-routed
boundary curves use the nested family
$A+B_1u+B_2u^2+H/\sqrt{\delta+1-u}$, allowing the degree-one subcase $B_2=0$,
where $u=(T-T_{\min,e})/(T_{\max,e}-T_{\min,e})$.  Routing uses only the
spread-support slope ratio (Appendix~\ref{app:protocol}); the sealed confirmation
package freezes $(\gamma,d,\delta)=(100,2,3\times10^{-4})$.
Coefficients denote background $A$, normalized slope $B_1$, curvature $B_2$,
and support-diagnosed boundary enhancement $H$---not microscopic constants.
DOI-grouped development reaches pooled $R^2=0.859300$ on 247 entities and
17,704 query rows, with
238/247 entities above $0.85$, median/p95/maximum NRMSE
$0.02958/0.21186/1.74495$, and 748 negative predictions.  The
inverse-square-root atom recurs in 5/5 folds, but only 1--3 training entities per
fold route to it; entity/DOI bootstrap lower bounds $0.82526/0.81641$ remain
below the endpoint.

\begin{table}[t]
  \centering
  \footnotesize
  \caption{Crystal-Cp development pooled physical $R^2$ with identical unseen-
  entity spread support.  CNP/no-q MLP are pointwise five-seed medians; other
  rows are deterministic.}
  \label{tab:crystal-development}
  \begin{tabular}{@{}lr@{}}
    \toprule
    Method & Physical $R^2\uparrow$ \\
    \midrule
    PCHIP & \textbf{0.958807} \\
    Linear interpolation & 0.952084 \\
    Nearest support & 0.919874 \\
    Support kNN & 0.898820 \\
    Boundary-stage expression & 0.859300 \\
    Shomate5 support re-q & 0.638646 \\
    FPCA & 0.632189 \\
    CNP & 0.346495 \\
    No-q temperature MLP & 0.095713 \\
    \bottomrule
  \end{tabular}
\end{table}

% ROLE: terminal development ablation -- separate a useful response chart from a failed learned prior.
The terminal $5$-fold $\times$ $5$-seed audit passes leakage checks and finds
more reproducible function than raw-latent geometry (mean Spearman $0.9230$ vs.
$0.6056$), yet both learned-bridge and GIRD gates fail.  Decoder-functional
spread $R^2$ is $0.4656$; in four-support inference, conditional GIRD median
NRMSE on the consensus prior-eligible subset is $0.0874$, worse than direct-
target dictionary $0.0497$ and $\lambda=0$ support OMP $0.0537$ on that subset.
Non-recurrent atoms, 0/300 passing sufficient-margin
certificates, and a failed gauge audit locate the problem in sparse-prior
identifiability (Appendix~\ref{app:protocol}).  The direct expression still
reaches $R^2=0.8593$, so this rejects the decoder prior rather than the
response-coordinate interface.

% ROLE: frozen pending branch -- reserve the only legal place for future evidence.
\paragraph{Single-use Crystal-Cp result.}
The 2015--2019 DOI/InChIKey-disjoint cohort is reserved under a single-use
contract.  The formal 25-cell matrix and all-development packages are terminal;
a target-blind seal binds 27 methods before any response is opened.  One
evaluation is authorized.  GIRD remains diagnostic, and the frozen expression
is unchanged even if it misses $R^2=0.85$.
